\documentclass{article}
\usepackage[T1]{fontenc}
\usepackage{lmodern}
\usepackage{graphicx}
\usepackage{amsmath,amssymb}
\usepackage[a4paper,margin=2cm]{geometry}
\usepackage[absolute,overlay]{textpos}
\setlength{\TPHorizModule}{1mm}
\setlength{\TPVertModule}{1mm}
\usepackage[indonesian]{babel}
\usepackage{hyperref}
\setlength{\emergencystretch}{2em}
% unit: Halaman 9

\begin{document}

% === Halaman 9 (llm) ===

\textbf{Contoh 2}: Misalkan bobot-bobot yang membentuk barisan \textit{superincreasing} adalah $\{2, 3, 6, 13, 27, 52\}$, dan diketahui bobot \textit{knapsack} $(M) = 70$. Kita akan mencari $b_1, b_2, ..., b_6$ sedemikian sehingga

\[ 70 = 2b_1 + 3b_2 + 6b_3 + 13b_4 + 27b_5 + 52b_6 \]

Caranya sebagai berikut:

\begin{enumerate}
    \item Bandingkan $70$ dengan bobot terbesar, yaitu $52$. Karena $52 \le 70$, maka $52$ dimasukkan ke dalam \textit{knapsack}. $\rightarrow b_6 = 1$
    \item Bobot total sekarang menjadi $70 - 52 = 18$. Bandingkan $18$ dengan bobot terbesar kedua, yaitu $27$. Karena $27 > 18$, maka $27$ tidak dimasukkan ke dalam \textit{knapsack}. $\rightarrow b_5 = 0$
    \item Bandingkan $18$ dengan bobot terbesar berikutnya, yaitu $13$. Karena $13 \le 18$, maka $13$ dimasukkan ke dalam \textit{knapsack}. $\rightarrow b_4 = 1$
\end{enumerate}

\end{document}
